\section{The Problem}

The reason to try this reversal is that the usual order has a hole. Start with matter, as physics does, and grant every success of that program, which is vast. You can explain the chemistry of a brain, the firing of its cells, the signals that lift a hand. You cannot explain why any of it is felt. A full account of the mechanism describes a system that processes information and steers its body in the dark, with nothing it is like to be it. That it is awake instead is left untouched. This is the hard problem.

The difficulty is structural. Physical description is description from the outside, what a thing does and how it relates to other things, and any such catalog could be satisfied by a system with no inside at all. Feeling is the one thing that is not a relation but a quality in itself, and a vocabulary of pure relation has no slot for it. Add mechanisms and the question only sharpens, because each mechanism is one more outside, and no amount of outside sums to an inside. The felt cannot be built from the unfelt, just as meaning cannot be built from marks that mean nothing, however many you stack.

The common moves deny the problem or quarantine it. Call the inner life an illusion, though to seem to feel pain is already to feel it. Call feeling an extra that appears when matter grows complex, though that names the mystery instead of solving it. Or set it aside as not yet ripe, which is how the science has progressed on everything else. None is a solution. When a question resists every answer of a given form, doubt the form. The form is matter first, feeling later. The doubt is whether the order can be reversed.

There is a long tradition on the other side of this. For thousands of years, and likely longer, the sages and contemplatives of many cultures have tried to set down what they found by looking inward, at the felt and conscious side of the world, and what they report is not easily carried in a few words, or even in many books. It comes wrapped in intricate and beautiful structure, in numbers and symmetries and geometries of the inner life, in maps drawn from long and disciplined experience.

I have never been willing to write all of that off as pure myth. In my own life I have caught a glimpse of what could be a deeper structure to reality, the same structure those traditions seem to circle, and this paper is an attempt to distill that glimpse into something more exact, a mathematical model, a map of a deep inner world that one can only really come to know through meditation, reflection, and patient introspection, approached with an open heart and an open mind. But a glimpse is not a proof, and an open mind with no hard edge will believe anything, so the whole of it is held to skepticism, to rationality and logic, and to a great deal of exploration and trial and error. What follows is what survived that test.

Modern science, for its part, has been extraordinarily good at the other half of the task, finding its slow way by trial and error to models that anyone can reproduce and that earn a public consensus, which is exactly the discipline the inner traditions lacked. But for all that power it has not yet arrived at the thing the sages were pointing toward, a single exact and shareable frame that holds the felt and conscious structure of the world, the beautiful and creative architecture they kept describing and hinting at across the ages. The two halves have stayed apart, the rigor on the one side and the inner vision on the other. This paper is an attempt to take a step toward closing that gap, to bring the discipline of the one to the subject matter of the other, and it is offered in that spirit, a step in what I believe is the right direction rather than a finished arrival.

This points to where the real problem lies. The inner side of reality has stayed outside science not because it is unreal but because we have lacked the measurable, observable handles that science runs on, no instrument turned the right way, no model in which the felt architecture became something one could compute and check. The inner depths needed better maps and better tools before they could be explored at all, and lacking them, science has had to leave a genuine part of reality unaddressed, the very part this paper sets out to render in intricate and exact detail. What has been missing, in a word, is the union of the two ways of knowing, the spiritual and esoteric and mystical on the one hand and the scientific and rational and mathematical on the other, each holding half of the picture and neither quite able to reach the other's half. This paper hopes to begin building that bridge, to give the inner architecture a structure precise enough to measure and to reason about, so that the depths the sages pointed to might at last be approached with the same tools that made the outer world knowable.

It may help to say where this particular picture came from, since it began not at a desk but in a kind of meditation, worked over a long while. The first question was the most basic one, what the absolute minimal basis of a universe would have to be, what survives when everything that can be stripped away is. In the end only one thing lasted, experience experiencing itself, and even the line between the act and the thing acted upon blurred until the two were one, a single change-that-is-an-object I first called by a made-up name, a tropon, before plainer words replaced it. The language simplified as I saw that short noun-verb words, four letters each, could carry the foundational pieces cleanly, the way a small vocabulary models a whole system in software. Experience became the vibe, and in its barest form it could be nothing more than three values, a pain, a peace, and a pleasure, which came to me one day as the smallest alphabet a feeling could have.

From there the picture grew. I asked what a vibe, a vibration, would actually do, how it would radiate and ripple and carry on, and if such a thing replicated, each one vibrating against its neighbors and stirring new ones in turn, what grew was an endless mesh of vibrations rippling off one another without end. That forced the harder question of where it could start, and the answer that kept returning was that it could not start from a true nothing, because a true nothing is unstable. To be nothing at all is already to be a distinction, a fact, a scrap of information, and a distinction makes a before and an after, which ripples, and the ripple ripples, on and on. From there the work turned concrete and almost austere. How would such vibes hold the least bit of their own experience, and read and write it to one another, with no high machinery, no hidden hands, nothing but the simplest discrete parts. The rest of this paper is the answer that search arrived at.
